\section{Method}
\label{sec:method}

Given only a few (typically 3-5) casually captured images of a specific subject, without any textual description, our objective is to generate new images of the subject with high detail fidelity and with variations guided by text prompts. Example variations include changing the subject location, changing subject properties such as color or shape, modifying the subject's pose, viewpoint, and other semantic modifications. We do not impose any restrictions on input image capture settings and the subject image can have varying contexts. We next provide some background on text-to-image diffusion models (Sec.~\ref{sec:t2i}), then present our fine-tuning technique to bind a unique identifier with a subject described in a few images (Sec.~\ref{sec:finetune}), and finally propose a class-specific prior-preservation loss that enables us to overcome language drift in our fine-tuned model (Sec.~\ref{sec:method_prior_pres}).

\subsection{Text-to-Image Diffusion Models}
\label{sec:t2i}
Diffusion models are probabilistic generative models that are trained to learn a data distribution by the gradual denoising of a variable sampled from a Gaussian distribution. Specifically, we are interested in a pre-trained text-to-image diffusion model $\hat\bx_\theta$ that, given an initial noise map $\bepsilon \sim \mathcal{N}(\bzero, \bI)$ and a conditioning vector $\bc=\Gamma(\bP)$ generated using a text encoder $\Gamma$ and a text prompt $\bP$, generates an image $\bx_{\text{gen}}=\hat\bx_\theta(\bepsilon, \bc)$. They are trained using a squared error loss to denoise a variably-noised image or latent code $\bz_t \coloneqq \alpha_t \bx + \sigma_t \bepsilon$ as follows:
\begin{equation}
\label{eq:diffusion}
    \Eb{\bx,\bc,\bepsilon,t}{w_t \|\hat\bx_\theta(\alpha_t \bx + \sigma_t \bepsilon, \bc) - \bx \|^2_2}
\end{equation}
where $\bx$ is the ground-truth image, $\bc$ is a conditioning vector (e.g., obtained from a text prompt), and $\alpha_t, \sigma_t, w_t$ are terms that control the noise schedule and sample quality, and are functions of the diffusion process time $t \sim \mathcal{U}([0, 1])$. 
A more detailed description is given in the supplementary material.  

\subsection{Personalization of Text-to-Image Models}
\label{sec:finetune}

% \begin{figure*}[t]
% \centering
% \includegraphics[trim={0 0 0 0},clip,width=\textwidth]{figures/system_fig.pdf}
% \caption{\textbf{Fine-tuning.} Given $\sim3-5$ images of a subject we fine tune a text-to-image diffusion in two steps: (a) fine tuning the low-resolution text-to-image model with the input images paired with a text prompt containing a unique identifier and the name of the class the subject belongs to (e.g., ``A [V] dog”), in parallel, we apply a class-specific prior preservation loss, which leverages the semantic prior that the model has on the class and encourages it to generate diverse instances belong to the subject's class using the class name in a text prompt (e.g., ``A dog”). (b) fine-tuning the super resolution components with pairs of low-resolution and high-resolution images taken from our input images set, which enables us to maintain high-fidelity to small details of the subject.}
% \label{fig:system_overview} 
% \end{figure*}
\begin{figure}[t]
\centering
\includegraphics[width=\columnwidth]{figures/main_scheme_1.pdf}
\caption{\textbf{Fine-tuning.} Given $\sim3-5$ images of a subject we fine-tune a text-to-image diffusion model with the input images paired with a text prompt containing a unique identifier and the name of the class the subject belongs to (e.g., ``A [V] dog”), in parallel, we apply a class-specific prior preservation loss, which leverages the semantic prior that the model has on the class and encourages it to generate diverse instances belong to the subject's class using the class name in a text prompt (e.g., ``A dog”). }
\label{fig:system_overview} 
\end{figure}

Our first task is to implant the subject instance into the output domain of the model such that we can query the model for varied novel images of the subject. One natural idea is to fine-tune the model using the few-shot dataset of the subject. Careful care had to be taken when fine-tuning generative models such as GANs in a few-shot scenario as it can cause overfitting and mode-collapse - as well as not capturing the target distribution sufficiently well. There has been research on techniques to avoid these pitfalls~\cite{Robb2020FewShotAO,ojha2021few,li2020few,mo2020freeze,Wang_2020_CVPR}, although, in contrast to our work, this line of work primarily seeks to generate images that resemble the target distribution but has no requirement of subject preservation. With regards to these pitfalls, we observe the peculiar finding that, given a careful fine-tuning setup using the diffusion loss from Eq~\ref{eq:diffusion}, large text-to-image diffusion models seem to excel at integrating new information into their domain without forgetting the prior or overfitting to a small set of training images. 

\paragraph{Designing Prompts for Few-Shot Personalization}
Our goal is to ``implant" a new (\textit{unique identifier}, subject) pair into the diffusion model's ``dictionary'' .
In order to bypass the overhead of writing detailed image descriptions for a given image set we opt for a simpler approach and label all input images of the subject ``a [identifier] [class noun]", where [identifier] is a unique identifier linked to the subject and [class noun] is a coarse class descriptor of the subject (e.g. cat, dog, watch, etc.). The class descriptor can be provided by the user or obtained using a classifier. We use a class descriptor in the sentence in order to tether the prior of the class to our unique subject and find that using a wrong class descriptor, or no class descriptor increases training time and language drift while decreasing performance. In essence, we seek to leverage the model's prior of the specific class and entangle it with the embedding of our subject's unique identifier so we can leverage the visual prior to generate new poses and articulations of the subject in different contexts.


\paragraph{Rare-token Identifiers}
We generally find existing English words (e.g. ``unique'', ``special'') suboptimal since the model has to learn to disentangle them from their original meaning and to re-entangle them to reference our subject. This motivates the need for an identifier that has a weak prior in both the language model and the diffusion model. A hazardous way of doing this is to select random characters in the English language and concatenate them to generate a rare identifier (e.g. ``xxy5syt00"). In reality, the tokenizer might tokenize each letter separately, and the prior for the diffusion model is strong for these letters. We often find that these tokens incur the similar weaknesses as using common English words.
Our approach is to find rare tokens in the vocabulary, and then invert these tokens into text space, in order to minimize the probability of the identifier having a strong prior. We perform a rare-token lookup in the vocabulary and obtain a sequence of rare token identifiers $f(\bVh)$, where $f$ is a tokenizer; a function that maps character sequences to tokens and $\bVh$ is the decoded text stemming from the tokens $f(\bVh)$. The sequence can be of variable length $k$, and find that relatively short sequences of $k=\{1, ..., 3\}$ work well. Then, by inverting the vocabulary using the de-tokenizer on $f(\bVh)$ we obtain a sequence of characters that define our unique identifier $\bVh$. For Imagen, we find that using uniform random sampling of tokens that correspond to 3 or fewer Unicode characters (without spaces) and using tokens in the T5-XXL tokenizer range of $\{5000, ..., 10000\}$ works well.

\subsection{Class-specific Prior Preservation Loss}
\label{sec:method_prior_pres}
In our experience, the best results for maximum subject fidelity are achieved by fine-tuning all layers of the model. This includes fine-tuning layers that are conditioned on the text embeddings, which gives rise to the problem of \textit{language drift}. Language drift has been an observed problem in language models~\cite{Lee2019CounteringLD,lu2020countering}, where a model that is pre-trained on a large text corpus and later fine-tuned for a specific task progressively loses syntactic and semantic knowledge of the language. To the best of our knowledge, we are the first to find a similar phenomenon affecting diffusion models, where to model slowly forgets how to generate subjects of the same class as the target subject.

Another problem is the possibility of \textit{reduced output diversity}. Text-to-image diffusion models naturally posses high amounts of output diversity. When fine-tuning on a small set of images we would like to be able to generate the subject in novel viewpoints, poses and articulations. Yet, there is a risk of reducing the amount of variability in the output poses and views of the subject (e.g. snapping to the few-shot views). We observe that this is often the case, especially when the model is trained for too long.

To mitigate the two aforementioned issues, we propose an autogenous class-specific prior preservation loss that encourages diversity and counters language drift. In essence, our method is to supervise the model with its \textit{own generated samples}, in order for it to retain the prior once the few-shot fine-tuning begins. This allows it to generate diverse images of the class prior, as well as retain knowledge about the class prior that it can use in conjunction with knowledge about the subject instance. Specifically, we generate data $\bx_\text{pr}=\bxh(\bz_{t_1}, \bc_\text{pr})$ by using the ancestral sampler on the frozen pre-trained diffusion model with random initial noise $\bz_{t_1} \sim \mathcal{N}(\bzero, \bI)$ and conditioning vector $\bc_\text{pr} \defeq \Gamma(f(\text{"a [class noun]"}))$. The loss becomes:
\begin{multline}
\label{eq:prior_pres}
    \E_{\bx,\bc,\bepsilon,\bepsilon^\prime,t}[w_t \|\hat\bx_\theta(\alpha_t \bx + \sigma_t \bepsilon, \bc) - \bx \|^2_2 + \\
    \lambda w_{t^\prime} \|\hat\bx_\theta(\alpha_{t^\prime} \bx_\text{pr} + \sigma_{t^\prime} \bepsilon^\prime, \bc_\text{pr}) - \bx_\text{pr} \|^2_2],
\end{multline}
where the second term is the prior-preservation term that supervises the model with its own generated images, and $\lambda$ controls for the relative weight of this term. Figure~\ref{fig:system_overview} illustrates the model fine-tuning with the class-generated samples and prior-preservation loss. 
Despite being simple, we find this prior-preservation loss is effective in encouraging output diversity and  in overcoming language-drift. We also find that we can train the model for more iterations without risking overfitting. We find that \(\sim \) 1000 iterations with $\lambda = 1$ and learning rate $10^{-5}$ for Imagen~\cite{saharia2022photorealistic} and $5 \times 10^{-6}$ for Stable Diffusion~\cite{rombach2022high}, and with a subject dataset size of 3-5 images is enough to achieve good results. During this process, $\sim 1000$ ``a [class noun]'' samples are generated - but less can be used. The training process takes about 5 minutes on one TPUv4 for Imagen, and 5 minutes on a NVIDIA A100 for Stable Diffusion.